\graphicspath{{\pathtoexos/Info/Figures/}}
%%%%%%%%%%%%%%%%%%%%%%%%%%%%%%%%%%%
\begin{exo}[PC][1][colle]{Gestion de table de données de F1}
    \Annale[Informatique PC]{Centrale}{2022}
    Depuis la création du championnat du monde de formule 1 en 1950, la fédération internationale de l'automobile (FIA) conserve l'ensemble des résultats des différents grands prix. Ces différentes informations sont stockées dans une base de données relationnelle dont un modèle physique simplifié est schématisé figure 6 .
    
    \begin{fig}[label]
        \centering
        \includegraphics[width=\columnwidth]{F1_SQL.png}
        \capt{Structure physique simplifiée de la base de données du championnat du monde de formule 1}
    \end{fig}
Cette base comporte les six tables listées ci-dessous avec la description de leurs colonnes:

\begin{itemize}
    \item la table Circuit de circuits utilisés pour le championnat du monde de formule 1
    \begin{itemize}
        \item ci\_id identifiant (entier arbitraire) du circuit (clef primaire)
        
        \item ci\_nom nom du circuit (généralement le nom de la ville ou de la région dans laquelle il se trouve)
        
        \item ci\_pays pays dans lequel se trouve le circuit
        
        \item la table GrandPrix des courses comptant pour le championnat du monde
        
        \item gp\_id identifiant (entier arbitraire) du grand prix (clef primaire)
        
        \item gp\_date date de la course
        
        \item ci\_id circuit sur lequel s'est déroulé la course
        
        \item gp\_nbt nombre de tours de circuit
        
    \end{itemize}
        \item la table Pilote des pilotes de formule 1
        
        \begin{itemize}
            \item pi\_id identifiant (entier arbitraire) du pilote (clef primaire)
            
            \item pi\_nom nom de famille du pilote
            
            \item pi\_prenom prénom du pilote
            
        \end{itemize}
            \item la table Constructeur des écuries de formule 1
            
            \begin{itemize}
                \item co\_id identifiant (entier arbitraire) de l'écurie (clef primaire)
                
                \item co\_nom nom de l'écurie
                
                \item co\_pays pays d'origine de l'écurie
                
                \item la table Participation des participants (couple pilote, constructeur) pour une course donnée
                
                \item pa\_id identifiant (entier arbitraire) de la participation (clef primaire)
                
                \item gp\_id la course
                
                \item pi\_id le pilote
                
                \item co\_id le constructeur
                
                \item pa\_cla le classement à l'arrivée de ce participant ou NULL s'il n'a pas terminé la course ou a été disqualifié - la table Tour des temps mis par un participant pour chaque tour de course, de clef primaire (pa\_id, to\_num)
                
                \item pa\_id identifiant (entier arbitraire) de la participation
                
                \item to\_num numéro du tour (de 1 à gp\_nbt)
                
                \item to\_temps le temps (en secondes) pour le tour correspondant
            \end{itemize}
                
            \end{itemize}

            \begin{quest}
                \item Pourquoi avoir utilisé la colonne pi\_id comme clef primaire de la table pilote et pas la colonne pi\_nom?
                
                \item Estimer le nombre de lignes de la table Tour.
                
                \item Écrire une requête SQL qui liste, par ordre chronologique, la date et le nom du circuit de toutes les courses qui se sont déroulées en France (ci\_pays = 'France').
                
                \item Écrire une requête SQL qui liste, pour chaque course de l'année 2021, le nom du circuit, le nom du pilote gagnant et son temps de course.
                
                \item Expliquer le résultat de la requête suivante
                
            \end{quest}
                %%%%%%%%%%%%%%%%%%%%%%%%%%%%%%%%%%%%%%%%%%%
% \begin{fig}[label]
% \centering
% \includegraphics[width=0.5\columnwidth]{2024_02_29_aef0b9cd50a734e3a3b4g-2}

% \subsection{Opérations et fonctions disponibles}
% \subsection{Fonctions}
% \begin{itemize}
%   \item $\mathrm{a} \% \mathrm{~b}$ calcule le reste de la division euclidienne de a par b, son signe est toujours celui de $\mathrm{b}$ $5 \% 3 \rightarrow 2,-5 \% 3 \rightarrow 1$.

%   \item isinstance $(0, t)$ teste si l'objet o est de type $t$ isinstance (-13, int) $\rightarrow$ True, isinstance ( $(1,2,3)$, tuple) $\rightarrow$ True.

%   \item range ( $\mathrm{n}:$ int) renvoie la séquence des $\mathrm{n}$ premiers entiers $(0 \rightarrow n-1)$ list (range(5)) $\rightarrow[0,1,2,3,4]$.

%   \item range( $\mathrm{d}:$ int, $\mathrm{f}:$ int, $\mathrm{s}: \mathrm{int})$ construit une séquence d'entiers espacés de $\mathrm{s}$ débutant à $\mathrm{d}$ et finissant avant $\mathrm{f}$ list (range $(0,-10,-2)) \rightarrow[0,-2,-4,-6,-8]$.

%   \item enumerate(s) itère sur la séquence $s$ en renvoyant, pour chaque élément de $s$, un couple formé de son indice et de l'élément considéré

% \end{itemize}

% list (enumerate $([3,1,4,1,5])) \rightarrow[(0,3),(1,1),(2,4),(3,1),(4,5)]$.

% $-\min (\mathrm{a}, \mathrm{b}, \ldots), \max (\mathrm{a}, \mathrm{b}, \ldots$.$) renvoie son plus petit (respectivement plus grand) argument$ $\min (2,4,1) \rightarrow 1, \max (" U n "$, "Deux") $\rightarrow$ "Un".

% \begin{itemize}
%   \item math.sqrt(x) calcule la racine carrée du nombre $x$.

%   \item round(n) arrondit le nombre $\mathrm{n}$ à l'entier le plus proche.

%   \item math.pi valeur approchée de la constante $\pi$.

% \end{itemize}

% \subsection{Opérations sur les listes}
% \begin{itemize}
%   \item len(u) donne le nombre d'éléments de la liste u
% \end{itemize}

% $\operatorname{len}([1,2,3]) \rightarrow 3$, len $([[1,2],[3,4]]) \rightarrow 2$.

% \begin{itemize}
%   \item u.count (e) renvoie le nombre d'éléments de la liste u égaux à e $[1,2,3,1,5]$. count (1) $\rightarrow 2$.

%   \item $\mathrm{u}+\mathrm{v}$ construit une liste constituée de la concaténation des listes $\mathrm{u}$ et $\mathrm{v}$ $[1,2]+[3,4,5] \rightarrow[1,2,3,4,5]$.

%   \item $\mathrm{n} * \mathrm{u}$ construit une liste constituée de la liste u concaténée $\mathrm{n}$ fois avec elle-même $3 *[1,2] \rightarrow[1,2,1,2,1,2]$.

% \end{itemize}

% $-u[i]=$ e remplace l'élément d'indice $i$ de la liste $u$ par e.

% \begin{itemize}
%   \item $u[i: j]=v$ remplace les éléments de la liste $u$ dont les indices sont compris dans l'intervalle $\llbracket i, j \llbracket$ par ceux de la séquence $\mathrm{v}$ (peut modifier la longueur de la liste $\mathrm{u}$ ).
% \end{itemize}

% — $\mathrm{u}$.append (e) ajoute l'élément e à la fin de la liste u (identique à $\mathrm{u}[\operatorname{len}(\mathrm{u}): \mathrm{z}=[\mathrm{e}]$ ). On suppose que cette opération a une complexité temporelle en $O(1)$.

% — $\mathrm{u}$. extend (v) ajoute les éléments de la séquence $\mathrm{v}$ à la fin de la liste u (identique à $\mathrm{u}[\operatorname{len}(\mathrm{u}):]=\mathrm{v}$. On suppose que cette opération a une complexité temporelle en $O(\operatorname{len}(\mathrm{v}))$.

% \begin{itemize}
%   \item u.insert (i, e) insère l'élément e à la position d'indice i dans la liste u (en décalant les éléments suivants) ; si $i>=l e n(u)$, e est ajouté en fin de liste (identique à $u[i: i]=[e]$ ).
% \end{itemize}

% — del $u[i]$ supprime de la liste u son élément d'indice $i$ (identique à $u[i: i+1]=[]$ ).

% \begin{itemize}
%   \item del $u[i: j]$ supprime de la liste $u$ tous ses éléments dont les indices sont compris dans l'intervalle $\llbracket i, j \llbracket$ (identique à $u[i: j]=[]$ ).

%   \item u.reverse() modifie la liste u en inversant l'ordre de ses éléments. On suppose que cette opération a une complexité temporelle en $O(\operatorname{len}(\mathrm{u}))$.

%   \item e in u teste si l'élément e apparait dans la liste u.

% \end{itemize}

% \subsection{Opérations sur les tableaux}
% \begin{itemize}
%   \item np.array(s) crée un nouveau tableau contenant les éléments de la séquence s. La forme de ce tableau est déduite du contenu de $\mathbf{s}$
% \end{itemize}

% np.array $([[1,2,3],[4,5,6]]) \rightarrow\left(\begin{array}{lll}1 & 2 & 3 \\ 4 & 5 & 6\end{array}\right)$.

% \begin{itemize}
%   \item a.ndim nombre de dimensions du tableau a.

%   \item a.shape tuple donnant la taille du tableau a pour chacune de ses dimensions.

%   \item a.size nombre total d'éléments du tableau a.

%   \item np.around (a) produit un tableau dont les éléments sont ceux du tableau a arrondis à l'entier le plus proche.

%   \item a @ b calcule le produit matriciel des deux tableaux a et b. Les dimensions de a et b doivent être compatibles. On suppose que la multiplication d'une matrice $n \times n$ et d'un vecteur a une complexité temporelle en $O\left(n^{2}\right)$.

% \end{itemize}

% \subsection{Module graphique turtle}
% Le module turtle, inspiré du langage Logo, simule une « tortue» robot qui se déplace sur l'écran. Cette tortue est munie d'un «crayon» qui, quand il est baissé, trace à l'écran la trajectoire de la tortue. Au début du programme, la tortue est située au centre de la fenêtre de visualisation, crayon baissé et orientée vers la droite de l'écran.

% \begin{itemize}
%   \item turtle.pendown(), turtle.penup() baisse ou lève le crayon.

%   \item turtle.right(a), turtle.left(a) fait pivoter la tortue sur place à droite ou à gauche de a degrés.

%   \item turtle.forward(n), turtle.back(n) fait avancer ou reculer la tortue de $n$ pixels (si $n$ est négatif, le mouvement est inversé) ; $n$ peut être un entier ou un nombre à virgule flottante.

%   \item turtle.circle $(r, a)$ fait parcourir par la tortue un arc de cercle de a degrés dont le centre est situé $r$ pixels à gauche de la tortue. Si r est positif, le centre est effectivement à gauche de la tortue et l'arc de cercle est parcouru dans le sens trigonométrique; si r est négatif, le centre est situé à droite de la tortue et l'arc de cercle est parcouru dans le sens horaire. Si a est positif, l'arc de cercle est parcouru en marche avant, la tortue tourne donc à gauche si $r$ est positif et à droite si $r$ est négatif. Si a est négatif, l'arc de cercle est parcouru en marche arrière, la tortue recule sur sa gauche si $r$ est positif et sur sa droite si $r$ est négatif.

%   \item turtle.home() ramène la tortue au centre de la fenêtre, orientée vers la droite de l'écran.

%   \item turtle.reset() efface la fenêtre et repositionne la tortue au centre, orientée vers la droite de l'écran.

% \end{itemize}

\subsection*{Formulaire : Fonctions SQL}
\begin{itemize}
  \item COUNT(*) fonction agrégative qui retourne le nombre de lignes chaque groupe.

  \item $\operatorname{COUNT}(\mathrm{e})$ fonction agrégative qui calcule le nombre de valeurs non nulles chaque groupe de lignes.

  \item $\operatorname{MIN}(e), \operatorname{MAX}(e)$ fonctions agrégatives qui calculent respectivement le minimum et le maximum de e pour chaque groupe de lignes; e est une expression de type numérique, chaine de caractère ou date, heure.

  \item $\operatorname{SUM}(e)$ fonction agrégative qui somme e pour chaque groupe de lignes; e est une expression de type numérique ou durée.

  \item EXTRACT (c FROM d) extrait la composante c de la date d, d peut être n'importe quelle expression de type date, timestamp, time ou interval et c peut prendre une des valeurs (liste non exhaustive) 'century', 'year ', 'quarter' (trimestre, 1 à 4), 'month' (mois, 1 à 12), 'day' (jour du mois, 1 à 31), 'dow' (jour de la semaine, 0 à 6,0 désignant dimanche), 'doy' (jour de l'année, 1 à 366).

\end{itemize}


\end{exo}
%%%%%%%%%%%%%%%%%%%%%%%%%%%%%%%%%%%